\documentclass[11pt]{article}
% ---------------------------------------------------------------------------
%  Genuine multipartite entanglement in benzene: canonical vs localized
%  orbitals, thermal behaviour, and validation.
%  Compile:  pdflatex benzene_gme_report.tex   (twice for TOC/refs)
% ---------------------------------------------------------------------------
\usepackage[a4paper,margin=2.3cm]{geometry}
\usepackage{amsmath,amssymb}
\usepackage{booktabs,array}
\usepackage{xcolor}
\usepackage{enumitem}
\usepackage[hidelinks]{hyperref}

\newcommand{\code}[1]{\texttt{\small #1}}
\newcommand{\CGME}{\ensuremath{C_{\mathrm{GME}}}}
\newcommand{\rdm}{\ensuremath{\Gamma^{(n)}}}

\title{\textbf{Genuine Multipartite Entanglement in a Molecular System:\\
       the Benzene $\pi$ System}}
\author{Khuram Shahzad}
\date{\today}

\begin{document}
\maketitle
\tableofcontents
\bigskip

% ===========================================================================
\section{Introduction}
% ===========================================================================
In previous work we computed the genuine multipartite entanglement (GME)
measure \CGME{} for the one-dimensional Hubbard ring at several system sizes
$L$, both for the ground state and for thermal (Gibbs) states. Here we extend
the same analysis to a real molecular system. Although the pipeline supports a
benchmark set of thirteen molecules, this report focuses on a single, physically
transparent case: \textbf{benzene} (C$_6$H$_6$) with the $\pi$-system active
space CAS$(6,6)$.

Benzene is the natural first molecular target because its six $\pi$ electrons
in six $\pi$ orbitals form the textbook molecular analogue of the
\emph{half-filled six-site Hubbard ring} studied before: one orbital per carbon
atom, arranged on a ring. This allows a direct, like-for-like comparison
between the model system and a real molecule. The report addresses three
questions:
\begin{enumerate}[nosep]
  \item How much genuine $n$-partite entanglement ($n=3,4,5$) does the benzene
        ground state carry?
  \item How does the answer depend on the \emph{orbital basis} --- canonical
        (delocalized) versus localized orbitals --- a question raised in the
        supervision discussion?
  \item How does the entanglement respond to temperature (Gibbs mixing of
        excited states)?
\end{enumerate}

% ===========================================================================
\section{Methodology}\label{sec:method}
% ===========================================================================

\subsection{System and active space}
Geometry and active space are taken from the curated VQE-benchmark dataset
(\code{molecules\_data.py}, basis \code{cc-pVDZ}). A restricted Hartree--Fock
calculation (PySCF) yields the molecular orbitals; the 18 lowest orbitals
(the $\sigma$ framework and cores) are frozen as a doubly-occupied core and the
next six orbitals --- the $\pi$ system --- are kept active:
CAS$(6,6)$, i.e.\ 6 electrons in 6 orbitals, Hilbert dimension
$\binom{6}{3}^2=400$ in the $S_z=0$ sector.

The active-space Hamiltonian is built with OpenFermion+PySCF as a number- and
spin-preserving sparse operator and diagonalised exactly (\code{eigsh}). Key
ground-state data:
\begin{center}
\begin{tabular}{@{}ll@{}}
\toprule
HF energy                    & $-230.718789$~Ha \\
CASCI ground-state energy $E_0$ & $-230.753558$~Ha \\
First excitation gap $E_1-E_0$  & $0.17530$~Ha \\
Spin                         & singlet ($\langle S^2\rangle=0$) \\
\bottomrule
\end{tabular}
\end{center}
These active spaces are exactly those optimised in the CUDA-Q VQE benchmark,
whose records show $|E_{\mathrm{VQE}}-E_{\mathrm{CASCI}}|\sim10^{-14}$~Ha ---
the exact diagonalisation used here produces the same correlated state that the
VQE targets.

\subsection{Two orbital bases: canonical and localized}\label{sec:bases}
The GME analysis is carried out in \emph{two} orbital bases.

\paragraph{Canonical (delocalized).}
The six active orbitals as delivered by RHF: the benzene $\pi$ stack
$a_{2u},\,e_{1g}\,(\times2),\,e_{2u}\,(\times2),\,b_{2g}$, each delocalized over
the whole ring and ordered by orbital energy. Orbital index $i$ refers to the
$i$-th active orbital in this energy ordering.

\paragraph{Pipek--Mezey localized.}
The same six-orbital space after a Pipek--Mezey (PM) unitary rotation. PM
maximises the atomic concentration of each orbital and preserves the
$\sigma/\pi$ separation. With a fourth-power PM functional and a multi-start
optimisation (the localizer is non-convex), the procedure converges to the
fully symmetric solution: \textbf{six equivalent atomic-like $p_z$ orbitals,
one per carbon}, each carrying $\approx42\%$ L\"owdin weight on its home carbon
(the remainder spills onto the neighbouring atoms --- the irreducible aromatic
delocalization). In this basis, orbital index $i$ means \emph{the $\pi$ orbital
of carbon $i$}: benzene becomes, orbital-wise, a six-site ring.

\paragraph{What changes and what does not.}
Localization is a unitary rotation \emph{within} the active space: the ground
state, its energy and the gap are strictly invariant (verified numerically:
identical $E_0$ in both bases). What changes is the \emph{tensor-product
partition} into orbitals --- and therefore the reduced density matrices and the
entanglement between them. \CGME{} is a property of a state \emph{and} a chosen
partition; comparing the two bases is thus physically meaningful, not
redundant: the canonical numbers measure entanglement between delocalized
molecular-orbital modes, the localized numbers between \emph{atomic sites},
directly comparable to the Hubbard ring.

A useful diagnostic of the basis change is the total double occupancy
$\langle\sum_j n_{j\uparrow}n_{j\downarrow}\rangle$ over the active orbitals:
\[
  \langle D\rangle_{\mathrm{canonical}} = 2.96 \qquad\longrightarrow\qquad
  \langle D\rangle_{\mathrm{localized}} = 1.33 .
\]
In the canonical basis the state is close to the RHF determinant (three doubly
occupied $\pi$ MOs, $\langle D\rangle\!\approx\!3$); in the localized (site)
basis the correlated $\pi$ electrons partially avoid double occupancy, exactly
as in a moderately interacting Hubbard ring.

\subsection{$n$-body reduced density matrices and their validation}
For every cluster of $n$ orbitals ($n=3,4,5$; all $\binom{6}{n}$ combinations)
the interleaved $n$-body spin RDM is assembled element-by-element,
$\rho[\mathrm{bra},\mathrm{ket}]
 =\langle\psi_0|\,\hat O_{\mathrm{bra,ket}}\,|\psi_0\rangle$,
where $\hat O$ is a product of one creation--annihilation pair per cluster
orbital (both acting on the \emph{same} orbital, so no inter-orbital hopping is
involved). The result is a $2^n\times2^n$ matrix over the spin configurations
of the cluster, normalised to unit trace. Its diagonal contains spin-configuration
populations; the off-diagonal elements are spin coherences whose signs encode
singlet ($-$) versus triplet ($+$) two-orbital couplings.

\paragraph{$S_z$ superselection check.}
Since $|\psi_0\rangle$ lies in the $S_z=0$ sector, every reduced state must be
block-diagonal in the cluster's total $S_z$: coherences between spin
configurations of different $S_z$ (e.g.\
$\langle\!\uparrow\uparrow\downarrow\!|\rho|\!\downarrow\downarrow\uparrow\!\rangle$)
must vanish identically. This was verified explicitly over \emph{all} stored
molecular RDMs: the largest cross-$S_z$ matrix element is $0$ to machine
precision. The RDM heatmap pages now include an $S_z$-ordered view in which
this block structure is manifest.

\paragraph{The $n=6$ (whole-active-space) case.}
The interleaved \emph{spin} RDM presumes each cluster orbital is singly
occupied. For $n=6$ in the canonical basis the probability that all six MOs are
simultaneously singly occupied is numerically zero
($\langle D\rangle\!\approx\!3$: bonding MOs are doubly occupied), so the
$n{=}6$ spin RDM is ill-defined there and is skipped by construction. In the
localized basis single occupancy is far more probable
($\langle D\rangle=1.33$), and the $n=6$ analysis becomes meaningful; it is
left as a next step.

\subsection{GME measure and thermal states}
For mixed RDMs, \CGME{} is evaluated as the convex roof
$C(\rho)=\min_{\{p_i,\psi_i\}}\sum_i p_i\,C(\psi_i)$, minimised with the GCG
solver over decompositions of increasing rank with several random restarts; the
pure-state formula is used when $\mathrm{Tr}\,\rho^2\approx1$. Analytic
lower/upper bounds are stored alongside for the audit's bracket check.

Thermal states are Gibbs mixtures over the $k=10$ lowest eigenstates,
\[
  \rdm(T)=\sum_{m} p_m(T)\,\rdm[\psi_m],\qquad
  p_m(T)=\frac{e^{-(E_m-E_0)/T}}{\sum_{m'}e^{-(E_{m'}-E_0)/T}},
\]
with $k_B=1$ and $T$ in Hartree, on the grid $T\in[0,0.3]$~Ha (16 points). The
relevant scale is the gap, $E_1-E_0=0.175$~Ha, so the window spans from the
pure ground state to strong mixing of the low-lying excited configurations.

% ===========================================================================
\section{Results}\label{sec:results}
% ===========================================================================

\subsection{Ground-state GME: canonical versus localized orbitals}
Table~\ref{tab:t0} summarises \CGME{} at $T=0$ for every body order, in both
bases (statistics over all $\binom{6}{n}$ clusters).

\begin{table}[h!]
\centering
\caption{Ground-state ($T=0$) GME of benzene CAS$(6,6)$ per body order $n$,
in the canonical (delocalized) and Pipek--Mezey localized bases.
``\#GME'' counts clusters with $\CGME>0.05$.}\label{tab:t0}
\begin{tabular}{@{}c c ccc c ccc@{}}
\toprule
 & & \multicolumn{3}{c}{\textbf{Canonical}} & \phantom{x} &
     \multicolumn{3}{c}{\textbf{Localized (PM)}} \\
\cmidrule{3-5}\cmidrule{7-9}
$n$ & \#clusters & $\max\CGME$ & mean & \#GME & &
      $\max\CGME$ & mean & \#GME \\
\midrule
3 & 20 & 0.074 & 0.009 & 2 & & 0.006 & 0.002 & 0 \\
4 & 15 & 0.108 & 0.018 & 1 & & 0.221 & 0.058 & 3 \\
5 &  6 & 0.000 & 0.000 & 0 & & 0.128 & \textbf{0.123} & \textbf{6} \\
\bottomrule
\end{tabular}
\end{table}

\noindent\textbf{Finding 1: the two bases give qualitatively different --- and
complementary --- pictures.} In the \emph{canonical} basis the entanglement is
modest and confined to small clusters (best $n{=}3$ cluster $0.074$, best
$n{=}4$ cluster $0.108$), and vanishes identically for all $n{=}5$ clusters.
In the \emph{localized} basis the pattern inverts: three-site clusters are
essentially separable, four-site clusters reach $0.221$, and \textbf{every}
five-site cluster is genuinely multipartite entangled, with a remarkably
uniform value ($\mathrm{mean}=0.123$, $\max=0.128$ --- the six clusters are
symmetry-equivalent up to the site removed).

\textbf{Finding 2: in the localized (atomic-site) basis, benzene behaves like
the half-filled Hubbard ring.} GME that \emph{grows} with cluster size, uniform
values across symmetry-equivalent clusters, and suppressed double occupancy
($\langle D\rangle=1.33$) are precisely the fingerprints observed in the $L=6$
Hubbard ground state at moderate coupling. This supports the physical picture
of the benzene $\pi$ system as an effective six-site ring of correlated
spins --- and confirms that the \emph{site-local} basis is the one in which the
multipartite structure of aromatic correlation is visible. The canonical basis,
being close to the RHF determinant ($\langle D\rangle\approx3$), hides most of
it.

\subsection{Thermal GME (canonical basis)}
Table~\ref{tab:thermal} shows the evolution of the strongest cluster per body
order under Gibbs mixing (canonical basis; localized thermal runs are queued as
a next step).

\begin{table}[h!]
\centering
\caption{Thermal GME of benzene CAS$(6,6)$ (canonical basis): the maximum
\CGME{} over clusters at the ends of the temperature window, and the observed
peak. $k=10$ eigenstates, $T$ in Hartree ($k_B{=}1$), gap $=0.175$~Ha.}
\label{tab:thermal}
\begin{tabular}{@{}c cccc l@{}}
\toprule
$n$ & $\max\CGME(0)$ & $\max\CGME(0.3)$ & peak & at $T$ & strongest cluster \\
\midrule
3 & 0.074 & 0.447 & 0.447 & 0.30 & $(2,3,4)$ \\
4 & 0.108 & \textbf{0.769} & \textbf{0.769} & 0.30 & $(1,2,3,4)$ \\
5 & 0.000 & 0.026 & 0.037 & 0.14 & $(0,1,2,3,5)$ \\
\bottomrule
\end{tabular}
\end{table}

\noindent\textbf{Finding 3: thermal mixing \emph{creates} multipartite
entanglement among the canonical orbitals.} Unlike the Hubbard ring at small
$t/U$ --- where heating monotonically destroys the ground-state GME --- the
benzene canonical-basis \CGME{} \emph{rises} steeply with temperature, from
$0.108$ to $0.769$ ($n{=}4$) across the window. The mechanism is the same as
the weak thermal enhancement seen for near-separable Hubbard states: the $T=0$
state is close to a single determinant in this basis (nearly separable), and
the Gibbs admixture of low-lying excited configurations (the window spans
$\sim\!1.7\times$ the gap) generates strong multi-orbital correlations. It is a
reminder that ``thermal death of entanglement'' presupposes an entangled ground
state; a near-product ground state can only gain. (Note $T=0.3$~Ha is an
enormous physical temperature $\sim\!10^5$~K; the Gibbs state is used here as a
controlled family of mixed states parameterised by the gap, not as an
experimentally reachable regime.)

\subsection{Numerical floor}
GCG runs with a bounded iteration budget, so near-separable RDMs return a small
positive residual ($\CGME\sim10^{-2}$) rather than exactly zero; the audit's
bracket check flags these as \textsc{warn}. All qualitative statements above
(which clusters are entangled, the canonical/localized inversion, the thermal
trends) are far above this floor.

% ===========================================================================
\section{Data products, auditing and reproducibility}\label{sec:saved}
% ===========================================================================
Every computed object is cached as a self-contained PKL, mirroring the Hubbard
pipeline:
\begin{itemize}[nosep]
  \item \code{mol\_gme\_runs/Benzene/molgme\_v1\_Benzene\_CAS66\_sites-<a-b-c>\_n<n>.pkl}
        --- per-cluster $T{=}0$ records (canonical); the localized runs carry a
        \code{\_locpm} tag. Each stores the full metadata (energies, gap, spin,
        double occupancy), the RDM (real/imaginary parts, spectrum, purity,
        entropy), all \CGME{} values with bounds, and the GCG decompositions.
  \item \code{mol\_thermal\_runs/Benzene/molthermal\_v1\_Benzene\_CAS66.pkl}
        --- the thermal curves $\CGME(T)$ and purities for every cluster.
  \item PDFs under \code{results/mol\_results/Benzene/}: RDM heatmaps (index
        order and $S_z$-ordered side by side), per-cluster GME pages, the
        $T{=}0$ bar chart and the thermal curves.
\end{itemize}
An automated audit (\code{main\_mol.py --audit}) checks completeness, the
lower/upper-bound bracket, RDM validity ($\mathrm{Tr}\rho=1$, hermiticity,
positivity), coverage, GME incidence and thermal-curve sanity, and exports tidy
CSV summaries. The present data pass all validity checks; the only
\textsc{warn} is the numerical floor discussed above.

% ===========================================================================
\section{Conclusions and next steps}
% ===========================================================================
For benzene CAS$(6,6)$ we find that (i) the ground state carries genuine
multipartite entanglement whose visibility depends strongly on the orbital
partition; (ii) in the Pipek--Mezey \emph{site-local} basis the $\pi$ system
reproduces the qualitative GME pattern of the half-filled six-site Hubbard
ring --- growth with cluster size, symmetry-uniform values, suppressed double
occupancy --- while in the canonical basis the state is nearly separable and
entanglement is instead \emph{generated} by thermal mixing; and (iii) all
reduced states rigorously respect the $S_z$ superselection structure, as
required.

Next steps: thermal analysis in the localized basis; the $n=6$
(whole-$\pi$-system) GME in the localized basis, where it is well defined;
quantitative comparison with the $L=6$ Hubbard ring at matched effective
coupling (via $\langle D\rangle$); and extension of the identical pipeline to
the remaining benchmark molecules on the HPC cluster.

\end{document}
